\section{Implementation}
\label{sec:impl}

\tool is written in Python and uses only the standard library, so the
competition container needs no extra packages. It has \NLocSpec lines of
code (without blank lines and comments): \NLocParser for Step~1,
\NLocDeps for Step~2, and \NLocValues for Steps 3 and~4. The shared data
types (operation, parameter, request, response) come from the team's common
interface file and are not counted.

\paragraph{Tests.} The tests have \NLocTests lines. A ladder of \NLevels
small specifications checks Step~1 one rule at a time, from a single
\texttt{GET} up to recursive schemas, \texttt{allOf} and base paths, against
exact expected output. A second check parses all \NSpecs real specifications
and counts the operations and the references left. A third check generates
30 requests per operation and mode for every specification and fails on a
crash or on a \mode{VALID} value that breaks its schema. A fourth check uses
a stand-in pool and asserts that real ids reach the path parameters.

\paragraph{Pitfalls.} Four problems showed up only on real specifications
and APIs.
First, the flight-search specification gives some strings a
\texttt{maxLength} of $2^{31}-1$. An early version of \mode{BOUNDARY} mode
tried to build a string of that length and ran out of memory. \tool now caps
strings at 5{,}000 characters and arrays at 50 items. Second, one
specification writes its patterns in Java syntax, such as \verb|\p{L}|,
which Python cannot parse. \tool then skips its pattern generator and keeps
the values from the other sources. Third, a first live evaluation never
created a hospital on gestao-hospital. Its specification lists an
\texttt{id} field in the body of \texttt{POST /v1/hospitais/} without marking
it read-only, so \tool sent one. The API then looks for a hospital with that
\texttt{id}, finds none, and answers 400; the same body without the
\texttt{id} returns 200
(\texttt{data/diagnosis.json}). \tool now leaves out an optional field named
\texttt{id} from the body of every \texttt{POST} request. Fourth, Step~2
removes a plural \texttt{s} the English way, which turns the Portuguese
collection \texttt{hospitais} into the key \texttt{hospitaiid}, while the path
parameter \texttt{hospital\_id} becomes \texttt{hospitalid}. The two never
matched, so no hospital id was ever reused. Two keys now also match when they
share their first six characters. All results in this report come from the
version with the last two changes.
